\section{Resultados}
Los resultados obtenidos tras la implementación del sistema reflejan avances 
significativos respecto a la situación inicial descrita en el Informe de Requerimientos 
[10]. Entre los logros más relevantes se destacan: 

1.  Centralización del registro de incidencias, lo que evita duplicidad de solicitudes y 

mejora la trazabilidad de cada caso. 

2.  Implementación de un sistema de roles que diferencia entre funcionarios, 
técnicos, administradores y superadministradores, garantizando un control 
personalizado del sistema. 

3.  Incorporación de notificaciones en tiempo real, que facilita la comunicación 

inmediata entre usuarios, técnicos y administradores. 

4.  Inclusión de un panel de métricas para administradores, que permite medir 

tiempos de atención, evaluar desempeño y generar reportes. 

5.  Desarrollo de un módulo de clasificación automática de tickets, optimizando la 

asignación de recursos. 

En términos de comparación con la literatura, estos resultados demuestran que Laravel 
cumple con las expectativas en materia de seguridad y rapidez [1], mientras que React 
garantiza una experiencia de usuario fluida y altamente interactiva [2]. Asimismo, la 
aplicación de Scrum aportó la flexibilidad necesaria para adaptarse a cambios en los 

requerimientos, lo cual coincide con lo planteado en estudios recientes sobre 
metodologías ágiles [9]. 